\section{Manutenzione straordinaria del campo da calcio in oratorio}

Il campo da calcio dell'oratorio di Maserada è un punto di incontro importante per i giovani della parrocchia. Attorno alla palla con l'arrivo della bella stagione si ritrovano a giocare squadre di parecchi ragazzini di tutte le età. Il divertimento però negli ultimi anni è stato compromesso dalle enormi pozzanghere d'acqua e fango che si formano dopo ogni pioggia proprio davanti alle due porte e che rimanendo anche per diversi giorni trasformano in una vera piscina olimpica il campo. Fenomeno questo che si è presentato anche durante il GrEst di inizio estate o in occasione dei raduni delle squadriglie di scout rendendo gli spazi non più utilizzabili. Insomma, un vero problema dovuto al fatto che il terreno non drena assolutamente l'acqua piovana:

\figureafiancoqui{porta_sud}{porta_nord}

Ebbene lo scorso sabato 28 febbraio una squadra di dieci volontari dopo settimane di preparativi in una sola giornata ha risolto la situazione! Sono stati eseguiti opportuni scavi per realizzare un sistema drenante composto da canali con ghiaia e tubo preforato il tutto avvolto da tessuto al fine di raccogliere l'acqua piovana davanti alle porte e convogliarla in pozzetti a fondo perduto a lato del campo.

Un lavoro di squadra magistrale eseguito con precisione. Ogni volontario ha messo a disposizione le sue competenze professionali e gli attrezzi da lavoro. È doveroso ringraziare lo sponsor che ha messo a disposizione l'escavatore e la Pro Loco di Maserada che ha coperto tutte le spese economiche. È stata una giornata molto faticosa ma accompagnata da sorrisi e battute scherzose durante i lavori.  

\figuramedia{0.7\textwidth}{scavo_nord}

Da notare che sono stati eseguiti anche altri lavori accessori tutto attorno al campo da calcio. La cosa molto bella è stato lo spirito di gruppo che ha rivisto assieme questi volontari dopo diversi anni. Don Federico ha salutato ad inizio lavori ed è tornato poi per il pranzo organizzato giù in oratorio dove una squisita pasta al ragù ha scaldato i cuori e del buon vino ha scaldato le guance e lo spirito goliardico.

\figuramedia{0.7\textwidth}{posa_tubi}

\figureafiancoqui{scavatore}{posa_tombino}

Dopo la pausa pranzo i lavori sono proseguiti attorno alla seconda porta da calcio e sono terminati verso sera. I volontari hanno sacrificato una intera giornata del loro tempo e lasciato a casa moglie e figli ma la soddisfazione di aver fatto qualcosa di utile per i giovani e la comunità li ha ampiamente ripagati della fatica. Ora altri lavori di manutenzione sono previsti fuori dall'oratorio perché ci sono gli alberi da potare, le reti da calcio da cambiare, le giostrine per i bimbi da sistemare le opere di giardinaggio da completare. L'obiettivo è ospitare in sicurezza all'aperto compleanni, feste, giochi, tornei sportivi, il GrEst, i Boy Scout e le varie occasioni pastorali. Accogliere ragazzi, ragazze e famiglie che hanno la possibilità di ritrovarsi in uno spazio pulito e sano. 

\firma{Massimiliano Guerra}

\figuramedia{0.62\textwidth}{pranzo}
